% ===========================================================================
%  sviluppatore/07-marketplace/02-lucchetto.tex
% ===========================================================================
\section{I due strati del permesso di pubblicare}\label{sez:lucchetto}

Chi può caricare un artefatto sul marketplace è deciso da due condizioni in
cascata, e confonderle è il modo tipico di sbagliare una modifica in
quest'area.

\figuraTikz{lucchetto-marketplace}
  {I due strati, nell'ordine in cui si valutano.}
  {lucchetto-marketplace}

\subsection{Primo strato: il lucchetto}

\begin{commento}{dashboard/community.py:30 $\cdot$ \cod{marketplace\_bloccato()}}
Il marketplace e' chiuso?

E' il PRIMO dei due strati che regolano la pubblicazione, e durante la
beta e' quello che conta: chiuso significa che nessuno carica niente — ne'
gratis ne' a pagamento — e nessuno compra. Prima al suo posto c'era un
interruttore che, spento, lasciava pubblicare gratis: non bastava, perche'
cio' che entra in vetrina ci resta e il filtro che decide cosa merita
starci (l'abbonamento) non era ancora in piedi.

Lo decide il sito, non l'applicazione: e' l'unico punto che chi usa il
programma non puo' modificare. In caso di dubbio risponde CHIUSO.
\end{commento}

\noindent
Tre proprietà, tutte deliberate:

\begin{itemize}
  \item \textbf{Lo decide il sito.} È l'unico permesso che chi usa il
        programma non può modificare agendo sulla propria macchina --- tutto
        il resto della dashboard è aperto (\vediSez{permessi}), questo no.
  \item \textbf{In caso di dubbio: chiuso.} Il \cod{except} finale restituisce
        \cod{True}.
  \item \textbf{Gli acquisti ne derivano.} \simbolo{vendite\_attive()} non è
        più un interruttore a sé: è \cod{not marketplace\_bloccato()}. Il
        commento dice perché resta come funzione: «due nomi per la stessa
        domanda sono due occasioni di rispondere diversamente».
\end{itemize}

\subsection{Secondo strato: l'abbonamento}

A marketplace aperto, serve uno dei tre livelli. I livelli sono quattro
contando chi non è abbonato:

\begin{codiceInterno}{python}{I quattro livelli}{dashboard/community.py:98}
# I QUATTRO LIVELLI, con i nomi che l'interfaccia mostra. Il livello 0 non e'
# un ripiego: e' un utente pieno, che sfoglia, scarica, compra e usa. Non
# mette soltanto roba in vetrina.
NOMI_LIVELLO = {0: "Free", 1: "Author", 2: "Studio", 3: "Desk"}
\end{codiceInterno}

\begin{center}
\begin{tabular}{@{}R{10mm}L{20mm}L{108mm}@{}}
\toprule
\textbf{Liv.} & \textbf{Nome} & \textbf{Che cosa può fare} \\
\midrule
\rigaTenue
\num{0} & Free & Tutto il sistema: sfoglia, scarica, compra, usa. \emph{Non}
  pubblica. \\
\num{1} & Author & In più: pubblicare sul marketplace. \\
\rigaTenue
\num{2} & Studio & Idem, prezzo diverso. \\
\num{3} & Desk & Idem, prezzo diverso. \\
\bottomrule
\end{tabular}
\end{center}

\begin{nota}{Perché oggi i tre piani fanno la stessa cosa}
È uno stato intermedio dichiarato, non una dimenticanza. Durante la beta i
tre livelli concedono solo la pubblicazione e differiscono per prezzo; la
decisione su che cosa aggiungere a Studio e Desk --- o se ridurli a uno solo
--- è rimandata all'uscita dalla beta. I prezzi attuali sono segnaposto e si
cambiano dal pannello di amministrazione senza toccare il codice.
\end{nota}

\subsection{Gli amministratori}

\begin{commento}{dashboard/community.py:111 $\cdot$ \cod{livello\_utente()}}
La categoria dell'utente: 0 se non abbonato, 1..3 secondo il piano.

Gli amministratori risultano al livello massimo, perche' possono fare
tutto cio' che fa il livello massimo e anche di piu'. Non sono un livello
vero — un livello scade, loro no — ma all'interfaccia serve un numero.
\end{commento}

\noindent
Un amministratore può inoltre installare un artefatto a pagamento senza
comprarlo (\rotta{POST}{/api/marketplace/<id>/installa-admin}) e,
simmetricamente, può scegliere di pagarlo comunque --- il che per un
amministratore equivale a una donazione all'autore, visto che non ne avrebbe
bisogno.

\begin{attenzione}{Gli amministratori non attraversano i percorsi d'errore}
È il motivo per cui il difetto del token (\vediSez{ponte}) è rimasto
invisibile per mesi: in beta gli unici a pubblicare erano gli
amministratori, che passano da un'altra strada. Chi collauda una modifica in
quest'area deve farlo con un utente \emph{non} amministratore, altrimenti
collauda il ramo sbagliato.
\end{attenzione}
